% from ega4-15-terminal-fr.tex
\oldpage[IV-3]{247}
on peut supposer $X'$ connexe (en remplaçant au besoin $X'$ par celle de ses composantes connexes contenant $Z_y$). Mais alors le même raisonnement qu'au début prouve qu'on a nécessairement $Z\subset X'$; comme $h(Y)$ est connexe et contient $x'$, il est nécessairement contenu dans $X'$, donc $x=h(y)\in X'\cap f^{-1}(y)=Z_y$. C.Q.F.D.
\end{proof}

\begin{corollary}[15.7.10]
\label{IV.15.7.10-fr}
Soient $Y$ un préschéma localement noethérien, $f:X\to Y$ un morphisme séparé de type fini, universellement submersif (15.7.8) et tel que toutes les fibres $X_y=f^{-1}(y)$ soient géométriquement connexes. Alors, si $y\in Y$ est tel que $X_y$ soit propre sur $\mathbf k(y)$, $f$ est propre au point $y$.
\end{corollary}

\begin{proof}
Compte tenu de (15.7.5), on est ramené au cas où $Y=\operatorname{Spec}(\mathcal O_y)$ puisque les hypothèses sont stables par changement de base. Mais dans ce cas il suffit d'appliquer (15.7.9) à $Z=X$.
\end{proof}

\begin{corollary}[15.7.11]
\label{IV.15.7.11-fr}
Soit $f:X\to Y$ un morphisme séparé de présentation finie; supposons qu'il existe un morphisme surjectif de présentation finie $g:Y'\to Y$, qui soit en outre propre ou plat, et tel que si l'on pose $X'=X\times_Y Y'$, il existe une $Y'$-section de $X'$ (ou encore un $Y$-morphisme $Y'\to X$ (\textbf{I}, 3.3.14)). Supposons enfin que les fibres $X_y=f^{-1}(y)$ soient géométriquement connexes. Alors l'ensemble $U$ des $y\in Y$ tels que $X_y$ soit propre sur $\mathbf k(y)$ est ouvert dans $Y$, et la restriction $f^{-1}(U)\to U$ de $f$ est propre.
\end{corollary}

\begin{proof}
La question étant locale sur $Y$, on peut supposer $Y=\operatorname{Spec}(A)$ affine. Appliquant (8.9.1) à $f$ et à $g$, on voit qu'il existe un sous-anneau noethérien $A_0$ de $A$ et deux morphismes de type fini $f_0:X_0\to Y_0=\operatorname{Spec}(A_0)$, $g_0:Y_0'\to Y_0$ tels que $X=X_0\times_{Y_0}Y$, $Y'=Y_0'\times_{Y_0}Y$, $f=(f_0)_{(Y)}$ et $g=(g_0)_{(Y)}$; en outre, utilisant (8.10.5, (v), (vi) et (xii)) et (11.2.7), on peut supposer $f_0$ séparé, $g_0$ surjectif et propre (resp. plat) si $g$ est propre (resp. plat); utilisant (8.8.2), on peut de plus supposer qu'il existe une $Y_0'$-section de $X_0$. D'autre part, utilisant (9.7.7), on peut appliquer (9.3.3) à la propriété d'être géométriquement connexe, et on peut donc supposer $A_0$ choisi de sorte que les $f_0^{-1}(y_0)$ soient géométriquement connexes pour tout $y_0\in Y_0$. Enfin, utilisant (2.7.1, (vii)), il revient au même de dire que $f^{-1}(y)$ est propre sur $\mathbf k(y)$ ou que $f_0^{-1}(y_0)$ est propre sur $\mathbf k(y_0)$, où $y_0$ est la projection de $y$ dans $Y_0$. Ces remarques montrent qu'on est ramené à démontrer le corollaire lorsque $Y$ est \emph{noethérien}. Notons maintenant le

\begin{lemma}[15.7.11.1]
\label{IV.15.7.11.1-fr}
Soient $f:X\to Y$ un morphisme, $g:Y'\to Y$ un morphisme surjectif, qui est en outre propre ou plat et quasi-compact. Posons $X'=X\times_Y Y'$, $f'=f_{(Y')}:X'\to Y'$. Alors, si $f'$ est submersif (resp. universellement submersif), il en est de même de $f$.
\end{lemma}

\begin{proof}
On peut se borner au cas où $f'$ est submersif. En effet, supposons le lemme prouvé dans ce cas, et montrons que si $f'$ est universellement submersif, $f$ l'est aussi. Soit donc $Y_1\to Y$ un morphisme quelconque, et posons $Y_1'=Y'\times_Y Y_1$; si $X_1'=X'\times_{Y'}Y_1'$ et $f_1'=(f')_{(Y_1')}:X_1'\to Y_1'$, $f_1'$ est submersif par hypothèse; de plus le morphisme $g_1:Y_1'\to Y_1$ est surjectif, et propre (resp. plat et quasi-compact) si $g$ l'est. Donc, si l'on pose $X_1=X\times_Y Y_1$, $f_1=f_{(Y_1)}$, il résulte de l'hypothèse et du fait que $X_1'=X_1\times_{Y_1}Y_1'$, que $f_1$ est submersif, donc $f$ universellement submersif. Supposons donc seulement que $f'$ soit submersif. Il est clair tout d'abord que $f$ est surjectif. Soit $E$ une partie de $Y$ telle que $f^{-1}(E)$ soit fermé dans $X$; alors, si $p:X'\to X$ est la projection canonique,
